% Example lab handout. Compile: latexmk -xelatex lab-example.tex
\documentclass[10pt,aspectratio=169,t]{beamer}
\usetheme[lab]{upbminimal}

\course{Application Development for Mobile Devices}
\decklabel{Lab 3}
\title{Widgets \& UI}
\subtitle{Composing layouts and handling gestures}
\author{Dinu-Ștefan Rusu}
\institute{FILS, Universitatea Politehnica București}
\date{Week 6}

\begin{document}

\titleframe

\begin{frame}[eyebrow=Today]{What this lab is for}
  \vskip 2mm
  \begin{grid}[2]
    \cell[\faThLarge]{Compose a screen}{Build a todo list from Column, Row, ListView and Card, without a single hard-coded pixel offset.}
    \cell[\faHandPointer]{Handle a gesture}{Swipe a row away with Dismissible and keep the list consistent.}
  \end{grid}
  \vskip 6mm
  \taskmeta{Time}{2 hours}{Submit}{Push to your project repo, branch \code{lab-3}, before next lab}
\end{frame}

\begin{frame}[fragile,eyebrow=Setup]{Where your code goes}
  \begin{steps}
    \item Open the project from Lab 2 and create a branch:
\begin{shell}
git checkout -b lab-3
\end{shell}
    \item Create \code{lib/lab3/todo\_screen.dart}. Everything in this lab lives there.
    \item Run the app once before changing anything. If it does not build, fix that first.
  \end{steps}
  \begin{hint}
    Hot reload keeps your list contents between edits. Hot restart drops them: use it after changing \code{initState}.
  \end{hint}
\end{frame}

\begin{frame}[eyebrow=Task I]{The screen, element by element}
  \begin{columns}[T]
    \begin{column}{0.55\textwidth}
      \begin{steps}
        \item An AppBar titled \emph{Todos}.
        \item A ListView.builder with one Card per todo.
        \item A Checkbox in each Card, bound to \code{done}.
        \item A FloatingActionButton that appends a todo.
        \item Extract the Card into \code{TodoTile}.
      \end{steps}
    \end{column}
    \begin{column}{0.41\textwidth}
      \kv{Done when}{}
      \begin{checklist}
        \item Adding a todo scrolls the list, and the list does not jump.
        \item Toggling a checkbox rebuilds only that tile.
        \item \code{TodoTile} takes its data through the constructor.
      \end{checklist}
    \end{column}
  \end{columns}
  \begin{takeaway}{Reusable component means constructor in, callbacks out.}
    A tile that reaches into global state is not reusable.
  \end{takeaway}
\end{frame}

\begin{frame}[fragile,eyebrow=Task II]{The anatomy of one swipeable row}
  \begin{columns}[T]
    \begin{column}{0.55\textwidth}
\begin{dart}
Dismissible(
  key: ValueKey(todo.id),          // never the index
  direction: DismissDirection.endToStart,
  onDismissed: (_) => onDelete(todo.id),
  background: Container(color: Colors.red),
  child: TodoTile(todo: todo, onToggle: onToggle),
)
\end{dart}
    \end{column}
    \begin{column}{0.41\textwidth}
      \begin{steps}
        \item Wrap each tile in a Dismissible.
        \item Remove the item from the list in \code{onDismissed}.
        \item Show a SnackBar with an \emph{Undo} action.
      \end{steps}
      \begin{warning}
        Using the list index as the key makes the wrong row disappear after a delete.
      \end{warning}
    \end{column}
  \end{columns}
\end{frame}

\begin{frame}[eyebrow=Troubleshooting]{Layout errors and their fixes}
  \vskip 2mm
  \trouble{RenderFlex overflowed by 42 pixels}{A Row or Column has more content than space. Wrap the growing child in Expanded or Flexible.}
  \trouble{Vertical viewport was given unbounded height}{A ListView inside a Column. Give it a fixed height, or wrap it in Expanded.}
  \trouble{setState() called after dispose()}{An async callback finished after the screen was popped. Check \code{mounted} before calling setState.}
  \trouble{A dismissed Dismissible widget is still part of the tree}{You did not remove the item from the list in onDismissed.}
\end{frame}

\closingframe{Push before you leave.}{Branch lab-3 · one commit per task · screenshots in the PR description}

\end{document}
